\section{Discussion}

\subsection{Summary of Achievement}

We have shown that protein-to-genome alignment on paralogs is vulnerable to two independent failure modes---cross-species identity ceiling and locus confusion---and that TPRS, an alignment method with annotation-derived priors, addresses both by expanding the candidate pool along the dimensions where single-tool approaches break, then selecting the best candidate through learned routing. TPRS raises paralog primary exon F1 from 0.17--0.40 (best single tool) to 0.57--0.79, recovering 71--90\% of the one-to-one performance ceiling.

\textbf{From broken to actionable.} Before TPRS, paralog alignment was effectively broken: a single-tool approach yielded primary F1 of 0.17--0.40, making automated paralog alignment unreliable. TPRS raises the floor to a level where paralog alignment becomes practically useful---the resulting spliced alignments reach quality levels that enable reliable downstream use. The residual gap to the one-to-one ceiling (8--13\% in mammals, 29\% in zebrafish) is attributable to ambiguous paralogs with near-identical coding sequences (mammals) and greater evolutionary divergence (zebrafish), and sets a ceiling for sequence-only approaches.

\subsection{Positioning: An Alignment Method with Annotation-Derived Priors}

TPRS is an alignment method whose output is a spliced alignment---the same output type produced by miniprot, Spaln3, or any other protein-to-genome alignment tool. It is not a de novo gene predictor, a genome annotation pipeline, or a gene model refinement system. It differs from conventional alignment tools in one specific way: it incorporates annotation-derived auxiliary information (species-matched proteins and curated orthology) as priors that expand the candidate space, while still producing spliced alignments evaluated on the same alignment metrics.

This design is motivated by the causal diagnosis, not by a preference for using annotations. The forced-locus experiment (Section~2.3) establishes that annotation-free approaches face a provable physical ceiling: without species-matched sequence information, the best C-stage tool achieves only 0.64--0.71 F1 on paralogs even when given the correct locus. This is not an engineering limitation---it is an information-theoretic one. The human query protein simply lacks the sequence motifs needed to resolve exon boundaries in a diverged target species' paralog. Introducing annotation-derived priors (species-matched protein queries, co-ortholog locus recovery) is the only way to break this ceiling, and TPRS demonstrates that these priors, combined with learned routing, make the problem tractable.

For clarity: TPRS does not compete with genome annotation pipelines (BRAKER, GALBA, GeMoMa). Those pipelines integrate RNA-seq, ab initio gene prediction, and other evidence sources to produce complete gene sets; they operate at a different pipeline stage than TPRS. A fair comparison against annotation pipelines would require integrating TPRS into each pipeline as the protein evidence source and re-running the full pipeline---an engineering effort beyond this paper's scope. TPRS competes with the alignment tools that annotation pipelines currently embed as their protein evidence source, and the benchmark (Table~\ref{tab:benchmark}) is against those alignment tools on alignment metrics.

\subsection{Design Principles and Generality}

\textbf{Modularity as a design strength.} TPRS is designed as a framework, not a fixed pipeline. Each stage has a clean interface: Phase~1 outputs candidate loci, Phase~2 adds recovered loci, Phase~3 produces alignments and C-stage features, and Phase~4 performs selection. Individual components can be independently improved: better A-stage tools can be added to Phase~1 without touching downstream phases; alternative C-stage tools can replace miniprot in Phase~3; the XGBoost router can be substituted with any regressor in Phase~4. This modularity also means that TPRS can benefit from future advances in any single component without restructuring.

\textbf{Diagnosis-driven design as a general approach.} The methodology of causally decomposing a failure mode before designing a solution is generalizable beyond protein-to-genome alignment. The pattern---identify independent root causes through controlled experiments, quantify their contributions, and map each to a specific architectural decision---applies to any bioinformatics pipeline where a composite failure masks multiple underlying problems. The candidate-generation + learned-selection architecture is a general recipe for problems where (1) the search space is large, (2) no single heuristic dominates, and (3) the correct answer is in the candidate pool but ambiguous to select.

\subsection{Limitations}

TPRS has several limitations that define its current scope. (i) Phase~2 co-ortholog recovery relies on Ensembl Compara orthology annotations, limiting direct applicability to model organisms with curated orthology data. For non-model species, reciprocal best BLAST hit provides an alternative, though with potentially lower recovery precision. (ii) The learned router (Phase~4) is currently trained per species; cross-species generalization has been tested only preliminarily (Exp~7 D7, Section~4.4) and shows measurable degradation on evolutionarily distant species pairs. Systematic cross-species router evaluation, including pooled multi-species training vs.\ per-species training, remains future work. Deploying TPRS on a new species currently requires training data with ground-truth CDS annotations from that species. (iii) Phase~1 depends on the availability and parameter tuning of four external tools. (iv) Phase~3 requires species-matched protein sequences, which presupposes at least partial knowledge of the target species' proteome. (v) The computational cost of Phase~3 (dual miniprot runs per candidate locus) scales linearly with the number of candidates; in practice, mouse takes $\sim$2--3 CPU-hours for the full pipeline. These limitations reinforce TPRS's positioning as an alignment method for species with existing resources, not a de novo alignment solution.

\subsection{Future Work}

Several directions merit exploration. Cross-species router generalization (training one model on all species jointly) could reduce the per-species training requirement. Incorporating protein structural information (e.g., AlphaFold-predicted structures) may help resolve the hardest paralog cases where sequence alone is ambiguous. Extending TPRS to plant genomes, where polyploidy creates even larger paralog families, would test the pipeline's scalability. Replacing Ensembl Compara with reciprocal-best-hit BLAST for Phase~2 would extend TPRS to non-model organisms. Finally, integrating TPRS into automated genome annotation workflows as a paralog-aware protein evidence source could bring the benefits of accurate paralog alignment to the broader annotation community.
